
Considere a recorrência linear definida como
\begin{align*}
  f(n) = \alpha f(n-1) + \beta f(n-2), \\
  f(0) = c, \\
  f(1) = d.
\end{align*}

Considerando o método de obtenção de fórmulas fechadas usando diagonalização, reescrevemos esse sistema como
\[
  \underbrace{\begin{bmatrix}
    f(n) \\
    f(n+1)
  \end{bmatrix}}_{f^{(n)}}
  =
  \underbrace{\begin{bmatrix}
    0 & 1 \\
    \beta & \alpha
  \end{bmatrix}}_{A}
  \underbrace{\begin{bmatrix}
    f(n-1) \\
    f(n)
  \end{bmatrix}}_{f^{(n-1)}},
  \quad
  \text{ com }
  f^{(0)}
  =
  \begin{bmatrix}
    c \\
    d
  \end{bmatrix}
\]
para $n \geq 2$.

Considerando que $f^{(n)} = A f^{(n-1)}$, temos que $f^{(n)} = A^{n} f^{(0)}$.

Começamos diagonalizando essa matriz $A$ calculando seus autovalores e autovetores. Para encontrar seus autovalores, pela definição de autovalor:
\begin{calcproof}[.]
  \start{Av = \lambda v}
  \step[\calciff]{Álgebra}{
    Av - \lambda I v = \bar{0}
  }
  \step[\calciff]{Distributividade}{
    (A - \lambda I) v = \bar{0}
  }
  \step[\calciff]{Teorema da Matriz Inversível}{
    \det{A - \lambda I} = 0
  }
\end{calcproof}

Com isso, considerando nossa matriz $A$, obtemos:
\begin{calcproof}[.]
  \start{\det{A - \lambda I} = 0}
  \step[\calciff]{Definição ($A$) e Álgebra}{
    \det{\begin{bmatrix}
        - \lambda & 1 \\
        \beta & \alpha - \lambda
      \end{bmatrix}} = 0
  }
  \step[\calciff]{Cálculo do determinante $2 \times 2$ e Álgebra}{
    \lambda^2 - \alpha \lambda - \beta = 0
  }
  \step[\calciff]{Bhaskara}{
    \lambda = \dfrac{\alpha \pm \sqrt{\alpha^2 + 4\beta}}{2}
  }
\end{calcproof}

Para calcular os autovetores, fazemos:
\begin{calcproof}[.]
  \start{(A - \lambda I)v = \bar{0}}
  \step[\calciff]{Definição ($A$) e Álgebra}{
    \begin{bmatrix}
      -\lambda & 1 \\
      \beta & \alpha - \lambda
    \end{bmatrix}
    \begin{bmatrix}
      v_1 \\
      v_2
    \end{bmatrix}
    =
    \begin{bmatrix}
      0 \\
      0
    \end{bmatrix}
  }
  \step[\calciff]{Considerando a equação de cima (a de baixo é redundante)}{
    v_2 - \lambda v_1 = 0
  }
  \step[\calciff]{Álgebra}{
    v = \begin{bmatrix}
      1 \\ \lambda
    \end{bmatrix}
  }
\end{calcproof}

Dessa forma, obtemos
\[
  A
  =
  \begin{bmatrix}
    1 & 1 \\
    \lambda_1 & \lambda_2
  \end{bmatrix}
  \begin{bmatrix}
    \lambda_1 & 0 \\
    0 & \lambda_2
  \end{bmatrix}
  \dfrac{1}{\lambda_2 - \lambda_1}
  \begin{bmatrix}
    \lambda_2 & -1 \\
    - \lambda_1 & 1
  \end{bmatrix},
\]
considerando que a inversa de $V$ foi obtida simplesmente pela fórmula da inversa da matriz $2 \times 2$.

Como $f^{(n)} = A^{n} f^{(0)}$, temos que $f^{(n)} = V D^{n} V^{-1} f^{(0)}$. Com isso, temos
\[
  \underbrace{\begin{bmatrix}
    f(n) \\
    f(n+1)
  \end{bmatrix}}_{f^{(n)}}
  =
  \underbrace{\begin{bmatrix}
    1 & 1 \\
    \lambda_1 & \lambda_2
  \end{bmatrix}}_{V}
  \underbrace{\begin{bmatrix}
    \lambda_1 & 0 \\
    0 & \lambda_2
  \end{bmatrix}^n}_{D^{n}}
  \dfrac{1}{\lambda_2 - \lambda_1}
  \underbrace{\begin{bmatrix}
    \lambda_2 & -1 \\
    - \lambda_1 & 1
  \end{bmatrix}}_{V^{-1}}
  \underbrace{\begin{bmatrix}
    c \\
    d
  \end{bmatrix}}_{f^{(0)}}.
\]

Sabendo que
\[
  \begin{bmatrix}
    \lambda_1 & 0 \\
    0 & \lambda_2
  \end{bmatrix}^n
  =
  \begin{bmatrix}
    \lambda_1^n & 0 \\
    0 & \lambda_2^n
  \end{bmatrix},
\]
podemos finalmente calcular $f^{(n)}$ simplesmente multiplicando as matrizes $V D V^{-1}$ e o vetor $f^{(0)}$. Olhando apenas para o primeiro componente do vetor $f^{(n)}$ (pois queremos a fórmula fechada de $f(n)$, e não de $f^{(n+1)}$), obtemos
\[
  f(n) = \dfrac{\lambda_1^n (c \lambda_2 - d) + \lambda_2^n (d - c \lambda_1)}{\lambda_2 - \lambda_1},
\]
obtendo assim uma fórmula fechada para $f(n)$ que não depende dos valores anteriores da sequência.

Vale verificar que, para o caso da função de fibonacci, onde temos $\lambda_1 = \varphi$, $\lambda_2 = \psi$, $c=0$ e $d=1$, substituindo na fórmula fechada geral para $f(n)$, obtemos
\[
  f(n) = \dfrac{\varphi^n (0 \psi - 1) + \psi^n (1 - 0 \varphi)}{\psi - \varphi},
\]
que é o mesmo que
\[
  f(n) = \dfrac{\varphi^n - \psi^n}{\varphi - \psi},
\]
que é conhecida como a \emph{fórmula de Binet}, que cálcula o $n$-ésimo termo da sequência de Fibonacci por meio de uma fórmula fechada, assim como queríamos.
